% =====================================================================
\section{MIS capacitor: C--V characterization}\label{sec:mis}
% =====================================================================

\subsection{Expected maximum capacitance (item 1)}\label{sec:mis-cmax}
We measured the \SI{200}{\micro m}\,$\times$\,\SI{200}{\micro m} capacitor
(n$^+$ poly-Si / SiO$_2$ / p-Si), so $A=\SI{4.0e-4}{cm^2}$. In accumulation the measured
capacitance is the oxide capacitance, $C_{max}=C_{ox}=\varepsilon_{ox}A/t_{ox}$. For the target
thickness $t_{ox}=\SIrange{45}{50}{nm}$ we expect
\[
C_{max}=\frac{3.9\times\SI{8.85e-14}{F/cm}\times\SI{4.0e-4}{cm^2}}{t_{ox}}
\approx\SIrange{27.6}{30.7}{pF}.
\]

\subsection{High-frequency C--V curve and oxide thickness (items 2--3)}\label{sec:mis-hf}
\begin{figure}[tbp]
\centering
\begin{subfigure}[b]{0.46\linewidth}\centering
  \includegraphics[width=\linewidth]{mdc_mis_HF_1MHz.jpg}\caption{}
\end{subfigure}\hfill
\begin{subfigure}[b]{0.46\linewidth}\centering
  \includegraphics[width=\linewidth]{mis_bands.pdf}\caption{}
\end{subfigure}
\caption{(a) C--V curve of the MIS capacitor at \SI{1}{MHz}, from inversion to accumulation
(MDC screenshot of a second sweep; the values in the text come from the data file of the first one).
(b) Band diagrams of the p-type MIS structure in the three regimes.}
\label{fig:mis-hf}
\end{figure}
The capacitor was measured at \SI{1}{MHz} from inversion to accumulation
(\SI{+10}{V}$\to$\SI{-10}{V}, step \SI{0.1}{V}), see \cref{fig:mis-hf}a. For a p-type substrate we
can identify three regimes, with the band diagrams of \cref{fig:mis-hf}b:
\begin{itemize}
\item \textbf{Accumulation} ($V_g<V_{FB}\approx\SI{-4.3}{V}$): the negative gate voltage attracts
holes to the surface (bands bend up). The charge is right at the interface, so $C\approx C_{ox}$.
\item \textbf{Depletion} ($V_{FB}<V_g<V_T$): holes are pushed away from the surface and a
depletion region of width $W$ appears. Its capacitance $C_{Si}=\varepsilon_{Si}A/W$ is in series with
$C_{ox}$ (Eq.~3), so $C$ decreases.
\item \textbf{Inversion} ($V_g>V_T\approx\SI{-3}{V}$): the bands bend down enough to create a layer
of electrons at the surface. At \SI{1}{MHz} these electrons are too slow to follow the AC signal,
$W$ stops growing and $C$ reaches its minimum value $C_{min}$.
\end{itemize}
The whole curve is shifted towards negative voltages (the device is already in inversion at
$V_g=0$), which shows that there is a lot of positive charge in the oxide (\cref{sec:mis-extraction}).



\textbf{Oxide thickness.} The maximum capacitance is $C_{max}=\SI{30.6}{pF}$ (at \SI{-10}{V}), so
\[
t_{ox}=\frac{\varepsilon_{ox}A}{C_{max}}\approx\SI{45}{nm},
\]
which is inside the target range. The capacitance is still increasing slightly at \SI{-10}{V}
(\SI{30.9}{pF} at \SI{-15}{V}), so this value may be slightly overestimated (by less than \SI{1}{nm}).

\textbf{Remark on inversion.} In theory, at high frequency $C$ should stay at $C_{min}$ in inversion.
Here it increases from \SI{10.4}{pF} to \SI{16.6}{pF} at \SI{+10}{V}. A possible explanation is that,
because of the large positive oxide charge, the surface around the gate is also inverted and can
supply electrons to the region under the gate, so part of the inversion layer can still respond.

\subsection{Low frequency, illumination and hysteresis (item 4)}\label{sec:mis-lf}
\begin{figure}[htbp]
\centering
\begin{subfigure}[b]{0.325\linewidth}\centering
  \includegraphics[width=\linewidth]{mdc_mis_LF_1kHz.jpg}\caption{\SI{1}{kHz}, dark}
\end{subfigure}\hfill
\begin{subfigure}[b]{0.325\linewidth}\centering
  \includegraphics[width=\linewidth]{mdc_mis_light.jpg}\caption{\SI{1}{MHz}, with light}
\end{subfigure}\hfill
\begin{subfigure}[b]{0.325\linewidth}\centering
  \includegraphics[width=\linewidth]{mdc_mis_hysteresis.jpg}\caption{\SI{1}{MHz}, $\pm\SI{15}{V}$ loop}
\end{subfigure}
\caption{MDC screenshots of the other measurements on the same capacitor.}
\label{fig:mis-other}
\end{figure}

\begin{figure}[htbp]
\centering
\begin{subfigure}[b]{0.325\linewidth}\centering
  \includegraphics[width=\linewidth]{mis_overlay.pdf}\caption{}
\end{subfigure}\hfill
\begin{subfigure}[b]{0.325\linewidth}\centering
  \includegraphics[width=\linewidth]{mis_invC2.pdf}\caption{}
\end{subfigure}\hfill
\begin{subfigure}[b]{0.325\linewidth}\centering
  \includegraphics[width=\linewidth]{mis_maserjian.pdf}\caption{}
\end{subfigure}
\caption{(a) Comparison of the C--V curves at \SI{1}{MHz} (dark and with light) and at
\SI{1}{kHz}. (b) $1/C^2$ method on the \SI{1}{MHz} curve. (c) Maserjian function
$C^{-3}dC/dV_g$ of the \SI{1}{MHz} curve.}
\label{fig:mis-analysis}
\end{figure}

\textbf{Low frequency (\SI{1}{kHz}).} In accumulation and depletion the curve is similar to the
\SI{1}{MHz} one, but in inversion the capacitance goes back up to $C_{ox}$
(\cref{fig:mis-analysis}a). At \SI{1}{kHz} the electrons of the inversion layer have enough time to
follow the AC signal, so the charge changes at the surface, like in accumulation, and $C\approx C_{ox}$.
This is the expected low-frequency behaviour. The small plateau around \SI{-5}{V} could be due to
interface states, which respond at \SI{1}{kHz} but not at \SI{1}{MHz}. The curve is noisier because the
measured AC current is much smaller at low frequency.

\textbf{With light (\SI{1}{MHz}).} Light creates electron--hole pairs in the silicon. The electrons
go to the surface, so the inversion layer is formed and renewed much faster. As a result, the
capacitance in inversion increases (\SI{19.2}{pF} instead of \SI{16.6}{pF} at \SI{+10}{V}) and the curve
moves towards the low-frequency one. Accumulation and depletion hardly change, because they involve
holes (majority carriers). This is why C--V measurements must be done in the dark.

\textbf{Hysteresis.} Sweeping from \SI{+15}{V} to \SI{-15}{V} and back, the two curves do not
overlap: the curve measured on the way back is shifted by about \SI{-0.66}{V}. This corresponds to
$\Delta N=C_{ox}|\Delta V|/(qA)\approx\SI{3e11}{cm^{-2}}$ charges that are trapped and released in the
oxide during the sweep. The shift towards negative voltages after the accumulation part means that
positive charge (holes) was trapped near the interface. We also observed that $V_{FB}$ changes slightly
from one measurement to the next (between about \SI{-4.3}{V} and \SI{-4.8}{V}), so the oxide charge is
not very stable.

\subsection{Extraction of $N_A$, $V_{FB}$ and $N_s$ (item 5)}\label{sec:mis-extraction}
We used the \SI{1}{MHz} curve with $C_{ox}=\SI{30.6}{pF}$. In each case $N_s$ is calculated from
Eq.~9, $N_s=C_{ox}(\Phi_{MS}-V_{FB})/(qA)$, with $\Phi_{MS}\approx\SI{-0.92}{V}$ for the n$^+$ poly-Si gate.

\textbf{C$_{min}$--C$_{max}$ method.} With $C_{min}=\SI{10.4}{pF}$, Eq.~3 gives the silicon capacitance
$C_{Si}=\SI{15.7}{pF}$, i.e.\ $W_{max}=\varepsilon_{Si}A/C_{Si}=\SI{0.26}{\micro m}$. Solving Eq.~5 (by
iteration, as $N_A$ also appears in the logarithm) gives $N_A=\SI{1.3e16}{cm^{-3}}$. Then the Debye
length is $L_D=\SI{36}{nm}$ (Eq.~7) and the flat-band capacitance $C_{FB}=\SI{24.2}{pF}$ (Eq.~8).
Reading the voltage where $C=C_{FB}$ on the curve gives $V_{FB}=\SI{-4.32}{V}$.

\textbf{$1/C^2$ method.} In depletion, $1/C^2$ is a straight line of slope $2/(q\varepsilon_{Si}N_AA^2)$
(\cref{fig:mis-analysis}b). Fitting the linear part (\SIrange{-3.7}{-3.3}{V}) gives
$N_A=\SI{1.28e16}{cm^{-3}}$, and the intersection of the line with $1/C_{ox}^2$ gives
$V_{FB}=\SI{-4.10}{V}$. The result depends a little on the chosen fitting range.

\textbf{Maserjian function (bonus).} The function $Y=C^{-3}dC/dV_g$ (\cref{fig:mis-analysis}c) should
be flat in depletion, with $|Y|=1/(q\varepsilon_{Si}N_AA^2)$. In our case the flat part is reduced to a
minimum, which gives $N_A=\SI{1.23e16}{cm^{-3}}$. The point where $Y$ reaches one third of its minimum
(accumulation side) gives $V_{FB}=\SI{-4.09}{V}$.

\begin{table}[H]
\centering\small
\caption{Parameters of the MIS capacitor obtained with the three methods, compared with the values
given by the MDC software.}
\label{tab:mis}
\begin{tabular}{lcccc}
\toprule
 & C$_{min}$--C$_{max}$ & $1/C^2$ & Maserjian & MDC software\\
\midrule
$N_A$ (\si{cm^{-3}}) & \num{1.32e16} & \num{1.28e16} & \num{1.23e16} & \num{1.35e16}\\
$V_{FB}$ (V) & $-4.32$ & $-4.10$ & $-4.09$ & $-4.33$\\
$N_s$ (\si{cm^{-2}}) & \num{1.6e12} & \num{1.5e12} & \num{1.5e12} & \num{1.6e12}\\
\bottomrule
\end{tabular}
\end{table}

The three methods agree well: $N_A\approx\SI{1.3e16}{cm^{-3}}$ and $V_{FB}$ between \SI{-4.1}{V} and
\SI{-4.3}{V}. The charge density $N_s\approx\SI{1.5e12}{cm^{-2}}$ is positive and much larger than for a
good quality thermal oxide (typically \SIrange{e10}{e11}{cm^{-2}}). This charge is responsible for the
large negative shift of the curve.

\subsection{Pros and cons of the methods (item 6)}\label{sec:mis-methods}
\begin{table}[H]
\centering\small
\caption{Comparison of the extraction methods.}\label{tab:methods}
\begin{tabularx}{\linewidth}{p{2.4cm}XX}
\toprule
Method & Advantages & Drawbacks\\
\midrule
C$_{min}$--C$_{max}$ & Simple: only two values of the curve are needed. &
Depends on a good value of $C_{min}$ (not well defined here, since $C$ rises in inversion);
assumes a uniform doping.\\
$1/C^2$ & Uses the whole depletion region; can show if the doping is not uniform. &
The result depends on the chosen fitting range; the linear region is short here.\\
Maserjian & No fitting range to choose; $N_A$ and $V_{FB}$ read directly on the plot. &
Needs a derivative, so it is sensitive to noise; the flat part is not always visible.\\
\bottomrule
\end{tabularx}
\end{table}
